% Двата модела на слушане, поставени един до друг.
%
% Рисунката носи основното твърдение на глава I: при конвенционалните сървъри изборът
% на протокол стои в слушащия сокет, а при coroute се взема след приемането на връзката.
% Затова лявата половина има три слушащи единици, а дясната две, и разликата в броя
% трябва да се вижда, преди читателят да е прочел един ред текст.
%
% Черно-бяла рисунка. Разграничението е по форма, по дебелина на линията и по сив тон:
% сокетите са по-тъмни (gray!25), обработчиците по-светли (gray!10), класификаторът е
% ромб (gray!40). Няма нито един цвят.
\begin{figure}[htbp]
\centering
\begin{tikzpicture}[
    % Вътрешните координати са в сантиметри и стигат до около 15,7 по хоризонтала.
    % Мащабът 0,9 ги свива до 14,1 cm, тоест под ширината на текста 16,0 cm.
    scale=0.9, transform shape,
    sock/.style={proc, minimum width=2.9cm, minimum height=0.9cm, fill=gray!25,
                 font=\fontsize{8}{9.5}\selectfont},
    hnd/.style={proc, minimum width=2.8cm, minimum height=0.9cm, fill=gray!10,
                font=\fontsize{8}{9.5}\selectfont},
    % Целите на класификатора са по-ниски от останалите кутии, за да не се слеят в
    % един блок, и еднакво широки, за да се четат като алтернативи на един и същ избор.
    tgt/.style={hnd, minimum width=3.0cm, minimum height=0.72cm},
    cfg/.style={note, align=center, text width=2.7cm,
                font=\fontsize{7}{8.5}\selectfont},
    cls/.style={decision, fill=gray!40, font=\fontsize{8}{9.5}\selectfont},
    tag/.style={note, font=\fontsize{7}{8.5}\selectfont, inner sep=1.5pt},
    badge/.style={box, fill=gray!40, minimum width=4.6cm, align=center,
                  font=\fontsize{8}{10}\selectfont},
    head/.style={align=center, font=\fontsize{9.5}{11.5}\selectfont\bfseries},
    sub/.style={note, align=center, text width=6.9cm,
                font=\fontsize{7.5}{9}\selectfont},
  ]

  % ================================================================ заглавия
  \node[head] at (3.65, 1.05) {Конвенционален сървър};
  \node[sub]  at (3.65, 0.42) {изборът на протокол е записан в сокета, \mbox{тоест} е направен
                               преди първия пристигнал байт};

  \node[head] at (12.15, 1.05) {\code{coroute}};
  \node[sub]  at (12.15, 0.42) {изборът на протокол е направен след \code{accept}, по
                                първия прочетен байт};

  % Разделителна линия между двете половини. Пунктир, за да не се чете като стрелка.
  \draw[densely dashed, gray!60] (7.90, 1.35) -- (7.90, -10.35);

  % ================================================================= вляво
  % Всеки протокол носи собствен сокет и собствена конфигурация. Пунктираната рамка
  % огражда двете заедно, защото те не могат да бъдат разделени: това е втората цена,
  % за която говори разделът за Caddy.
  \node[sock] (s1) at (1.75, -1.15) {TCP, порт 80\\ \code{listen 80}};
  \node[cfg]  (c1) at (1.75, -2.25) {без сертификат\\ правило TCP 80};
  \node[grp, fit=(s1)(c1)] (g1) {};

  \node[sock] (s2) at (1.75, -4.00) {TCP, порт 443\\ \code{listen 443 ssl}};
  \node[cfg]  (c2) at (1.75, -5.10) {\code{ssl\_certificate}\\ правило TCP 443};
  \node[grp, fit=(s2)(c2)] (g2) {};

  \node[sock] (s3) at (1.75, -6.85) {UDP, порт 443\\ \code{listen 443 quic}};
  \node[cfg]  (c3) at (1.75, -7.95) {\code{ssl\_certificate}\\ правило UDP 443};
  \node[grp, fit=(s3)(c3)] (g3) {};

  \node[hnd] (h1) at (5.85, -1.61) {HTTP/1.1};
  \node[hnd] (h2) at (5.85, -4.46) {TLS, после\\ HTTP/1.1};
  \node[hnd] (h3) at (5.85, -7.31) {QUIC, HTTP/3};

  \draw[flow] (g1.east) -- (h1.west);
  \draw[flow] (g2.east) -- (h2.west);
  \draw[flow] (g3.east) -- (h3.west);

  % ================================================================= вдясно
  % Конфигурацията е обща и стои над сокетите, за да могат стрелките да излизат надолу
  % без да я пресичат.
  \node[note, align=center, text width=6.5cm, font=\fontsize{7}{8.5}\selectfont]
        (cr) at (12.15, -0.85)
        {обща конфигурация: един сертификат,\\ едно правило в защитната стена};
  % Един порт, не два. Един слушащ сокет се свързва с един номер на порт, така че
  % надпис „портове 80 и 443“ тук би означавал два сокета и следователно три
  % дескриптора, което противоречи както на значката отдолу, така и на предложението
  % L = W * |T| от раздел 2.3. Ползата е един изложен порт, а не наполовина по-малко
  % портове, и раздел 3.1 казва точно това.
  \node[sock] (st) at (10.30, -2.20) {TCP сокет\\ порт 443};
  \node[sock] (su) at (14.10, -2.20) {UDP сокет\\ порт 443};
  \node[grp, fit=(cr)(st)(su)] (gr) {};

  % Ромбът е единствената фигура с този контур в рисунката. Той е мястото, на което се
  % взема решението, и стои под приемането на връзката, а не над него.
  \node[cls] (cl) at (10.30, -4.30) {класификация по\\ първия октет};
  \draw[flow] (st.south) -- (cl.north)
        node[tag, midway, anchor=west, xshift=3pt] {\code{accept}};

  % Изходите на класификатора висят на обща вертикална линия, за да се четат като
  % алтернативи на един и същ избор, а не като последователни стъпки.
  \draw[thick] (cl.south) -- (10.30, -7.80);

  \node[tgt] (t1) at (13.50, -6.00) {HTTP/1.1};
  \node[tgt] (t2) at (13.50, -6.90) {HTTP/2, \code{preface}};
  \node[tgt] (t3) at (13.50, -7.80) {TLS ръкостискане};

  \draw[flow] (10.30, -6.00) -- (t1.west)
        node[tag, midway, above, xshift=-3pt] {ASCII метод};
  \draw[flow] (10.30, -6.90) -- (t2.west)
        node[tag, midway, above] {\code{PRI *}};
  \draw[flow] (10.30, -7.80) -- (t3.west)
        node[tag, midway, above] {\code{0x16}};

  \node[tgt] (al) at (13.50, -8.95) {ALPN избира\\ \code{h2} или \code{http/1.1}};
  \draw[flow] (t3.south) -- (al.north);

  \node[hnd] (qh) at (14.10, -4.30) {QUIC, HTTP/3};
  \draw[flow] (su.south) -- (qh.north);

  % ================================================================= броене
  % Числото е тук, защото то е сравнимата величина. Умножението по работните нишки не
  % изчезва при нито един от двата модела, променя се само какво се умножава.
  \node[badge] at (3.65,  -9.95) {\textbf{3} слушащи дескриптора\\
                                  $\times$ броя на работните нишки};
  \node[badge] at (12.15, -9.95) {\textbf{2} слушащи дескриптора\\
                                  $\times$ броя на работните нишки};

  \node[note, align=center, text width=13.0cm, font=\fontsize{8}{10}\selectfont]
        at (7.90, -10.85)
        {Умножението по броя на работните нишки остава и в двата случая.\\
         Променя се само числото, което се умножава.};

\end{tikzpicture}
\caption{Двата модела на слушане. Вляво: конвенционален сървър, при който всеки протокол
получава собствен слушащ сокет, а пунктираната рамка огражда сокета заедно с неразделната
му конфигурация. Изборът е направен преди пристигането на първия байт. Вдясно:
\code{coroute}, при който един TCP сокет и един UDP сокет са достатъчни, а протоколът се
разпознава по първия октет след \code{accept}, като при TLS изборът между \code{h2} и
\code{http/1.1} се доуточнява от ALPN.}
\label{fig:listener_models}
\end{figure}
